\documentclass[11pt]{article}
\usepackage{fontspec,amsmath,amssymb,amsthm,hyperref}
\usepackage[margin=1in]{geometry}
\setlength{\emergencystretch}{3em}
\title{Derived strong convexity of the actual marginal potential}
\author{LangevinLean mathematical exposition}
\date{}
\begin{document}\maketitle
\noindent Source module: \texttt{IPLA/StrongLogConcavity.lean}. Compile this proof document with XeLaTeX or LuaLaTeX. The accompanying HTML pairs the mathematical argument with the exact Lean statements and records the scope of the formalization.


Let \(M(\theta)=\int e^{-U(\theta,x)}dx\) and \(\Phi(\theta)=-\log M(\theta)\). The proved finite-dimensional Prékopa–Leindler theorem shows that \(\Phi\) retains the original joint strong-convexity constant \(m\).


\section{Define a legitimate logarithmic marginal potential}


The real marginal integral \(M\) is defined from the original slice; the earlier domination theorem makes it finite and strictly positive. Thus \(\Phi=-\log M\) is an ordinary real-valued continuous function: continuity follows by composing the continuous marginal with logarithm away from zero. If \(\theta_*\) maximizes the extended marginal, its finite real representation and monotonicity of logarithm imply \(\Phi(\theta_*)\le\Phi(\theta)\) for every \(\theta\). No logarithm of zero or infinite partition is used.




\textbf{Correspondence with the displayed declarations.}


\begin{itemize}

\item \texttt{marginalIntegral}: The real latent partition. Positivity and finiteness were proved from the joint assumptions in \texttt{MarginalConcentration.lean}.
\item \texttt{marginalPotential}: The actual marginal negative log-likelihood, with the observations fixed.
\item \texttt{marginalMaximizer\_minimizes\_potential}: The statistical marginal maximizer minimizes the actual negative log-likelihood. All real conversions use proved positive finite partitions.
\end{itemize}


\noindent Lean declarations: \texttt{marginalIntegral}, \texttt{marginalPotential}, \texttt{marginalIntegral\_pos}, \texttt{marginalPotential\_continuous}, \texttt{marginalMaximizer\_minimizes\_potential}.

\section{Prepare Gaussian weights in product coordinates}


Define \(f_\theta(x)=e^{-U(\theta,x)}\) in ordinary finite-product coordinates, using the volume-preserving Euclidean identification for the latent block. The joint Gaussian domination gives \(f_\theta(x)\le Ce^{-m\|x\|^2/4}\), and continuity plus strict positivity make it a Gaussian weight. Its integral equals \(M(\theta)\) by the volume-preserving change of coordinates. Finally pairing preserves subtraction and affine combinations coordinatewise:

\[\operatorname{pair}(\theta,x)-\operatorname{pair}(\theta',y)=\operatorname{pair}(\theta-\theta',x-y),\]\[\operatorname{pair}(a\theta+b\theta',ax+by)=a\operatorname{pair}(\theta,x)+b\operatorname{pair}(\theta',y).\]


\textbf{Correspondence with the displayed declarations.}


\begin{itemize}

\item \texttt{latentWeight}: A latent slice written in ordinary product coordinates for applying the finite-dimensional Prékopa--Leindler theorem.
\end{itemize}


\noindent Lean declarations: \texttt{latentWeight}, \texttt{latentWeight\_gaussian}, \texttt{integral\_latentWeight}, \texttt{pair\_sub}, \texttt{pair\_affine}.

\section{Apply Prékopa with the parameter convexity penalty}


Let \(a,b\ge0\), \(a+b=1\), \(\theta_r=a\theta+b\theta'\), and \(c=(mab/2)\|\theta-\theta'\|^2\). Strong convexity of \(U\), the affine pair identity and norm splitting imply

\[aU(\theta,x)+bU(\theta',y)\ge U(\theta_r,ax+by)+c,\]

after discarding the nonnegative latent penalty \((mab/2)\|x-y\|^2\). Exponentiate and apply Prékopa–Leindler with third weight \(e^{-c}f_{\theta_r}\), which remains a Gaussian weight. This gives \(M(\theta)^aM(\theta')^b\le e^{-c}M(\theta_r)\). All factors are positive, so take logarithms and negate to conclude

\[\Phi(a\theta+b\theta')\le a\Phi(\theta)+b\Phi(\theta')-\frac{mab}{2}\|\theta-\theta'\|^2.\]

Thus marginal strong convexity is a theorem of the joint assumptions.




For \(\theta_r=a\theta+b\vartheta\), \(a+b=1\), joint strong convexity yields \[U(\theta_r,ax+by)\le aU(\theta,x)+bU(\vartheta,y)-\tfrac{mab}{2}\|\theta-\vartheta\|^2,\] after dropping the additional nonnegative latent deficit. Thus the latent weights satisfy the Prékopa hypothesis with the multiplier \(e^{-c}\), \(c=mab\|\theta-\vartheta\|^2/2\), on the interpolated weight. Integrating gives \(M(\theta)^aM(\vartheta)^b\le e^{-c}M(\theta_r)\).



Every marginal integral is already finite and positive, so logarithms preserve the inequality. Negating gives \(F(\theta_r)\le aF(\theta)+bF(\vartheta)-c\), exactly \(m\)-strong convexity of \(F=-\log M\). The code's \texttt{hpoint} is the weight inequality, \texttt{hpl} its Prékopa integral consequence, and \texttt{hlog} its logarithmic conversion. Strict convexity for \(m>0\) then proves uniqueness of the previously constructed maximizer.




\textbf{Correspondence with the displayed declarations.}


\begin{itemize}

\item \texttt{marginalPotential\_strongConvex}: Integrating out the latent variables preserves precisely the original strong-convexity constant. This is the derived Prékopa theorem for IPLA.
\end{itemize}


\noindent Lean declarations: \texttt{marginalPotential\_strongConvex}.

\section{The marginal maximizer exists uniquely}


Because \(m>0\), the proved \(m\)-strong convexity makes \(\Phi\) strictly convex. Two marginal maximizers are both global minimizers of \(\Phi\); if distinct, the midpoint strict-convexity inequality would lie strictly below their common minimum, a contradiction. Hence they coincide. The earlier continuity-and-tail argument already supplied existence of a maximizer. Combining existence with this uniqueness yields \(\exists!\theta_*\) satisfying the original statistical maximizer predicate.




\textbf{Correspondence with the displayed declarations.}


\begin{itemize}

\item \texttt{marginalMaximizer\_unique}: Joint strong convexity implies uniqueness of the marginal maximizer.
\item \texttt{existsUnique\_marginalMaximizer}: The marginal maximizer exists uniquely under exactly the original joint potential assumptions.
\end{itemize}


\noindent Lean declarations: \texttt{marginalMaximizer\_unique}, \texttt{existsUnique\_marginalMaximizer}.

\end{document}
